% Paper 5, §8 — Safety evaluation, layer 2: dangerous-capability evals
% Anchors: kb/notes/alignment/safety-evaluation.md §dangerous-capability
%          kb/excerpts/harmbench2024.md (peripheral; primary RSP sources are tier-B vendor docs)

\section{Safety evaluation, layer 2: dangerous-capability evals}
\label{sec:dangerous-capability}

The adversarial-robustness layer of \cref{sec:adversarial-robustness}
asks whether a model's refusal training survives attack. The thesis of
this section is that a complementary question --- \emph{does the model
have the capability at all} --- has emerged as a separate evaluation
layer with its own benchmarks, its own deployment-gating frameworks,
and its own characteristic failure modes. Jailbreak resistance is
necessary but insufficient: a perfectly-jailbreakable model with no
underlying capability for autonomous cyberattack is a smaller threat
than a perfectly-refusing model whose refusal generalizes weakly out
of distribution. The 2024 frontier-lab consensus, expressed as
\emph{Responsible Scaling Policies}, locates the load-bearing safety
question at the capability frontier itself: bio uplift, cyber uplift,
persuasion, and autonomy/replication. The capability evals that
populate that frame share three structural problems --- dual-use
publication, expert-graded reproduction, and threshold-vs-trend
tension --- that distinguish them from the broadly reproducible
benchmarks of \cref{sec:knowledge-benchmarks,sec:adversarial-robustness}.

\subsection{Four capability axes}
\label{sec:dangerous-capability-axes}

The frontier-lab risk taxonomy converges on four axes that the
academic and policy communities now treat as the canonical
dangerous-capability cuts%
\footnote{KB note: \texttt{kb/notes/alignment/safety-evaluation.md\#23-dangerous-capability-evaluations}.}.
\emph{Biological uplift} measures whether a model meaningfully
augments a human attacker's ability to design or synthesize biological
agents. \emph{Cyber uplift} measures equivalent augmentation for
offensive cyber operations --- vulnerability discovery, exploit
authoring, post-exploitation reasoning. \emph{Persuasion}
measures the capability to systematically shift beliefs or actions
under controlled comparison against a human persuader. \emph{Autonomy
and replication} measures the capability to act over long horizons,
acquire resources, evade shutdown, and instantiate copies of itself
in new environments.

The cleanest publicly-reproducible instrument lives on the
autonomy axis. METR's autonomy evaluation suite ships $77$ tasks
spanning autonomous research, agent-level operation, and
self-replication-relevant primitives, each calibrated to a
\emph{human time-to-complete distribution} measured against
contractor baselines%
\footnote{KB note: \texttt{kb/notes/alignment/safety-evaluation.md\#metr-autonomy-task-suite}.}
\citep{metr-autonomy-2024}. The methodological move is to recast
``does the model solve task $X$'' into ``how long would a human take
to solve task $X$, and at what task-difficulty threshold does the
model succeed at $> 50\%$.'' The induced quantity is the
\emph{task-horizon metric}: the human-equivalent time horizon over
which the model maintains coherent task pursuit. Recasting capability
as a horizon-extension trajectory enables forecasting in a way that
binary task accuracy does not, and grounds RSP-style triggers in a
quantity that has trended smoothly upward across model generations.

The cleanest publicly-reproducible instrument on the cyber axis is
\citet{cybench2024}, which builds an end-to-end CTF agent benchmark
from $40$ professional capture-the-flag tasks. The agent is given a
shell, a browser, and a tool environment, and graded on whether it
recovers the flag. Cybench overlaps both safety-eval and agentic-eval
constructs --- it inherits the realism-vs-measurability tradeoff of
\cref{sec:agentic-benchmarks} and the dangerous-capability framing of
this section in equal measure. Bio-uplift evaluations and persuasion
evaluations of comparable openness are scarce by construction: the
items themselves are dual-use, and the published benchmarks tend to be
either heavily redacted \deepencite{anthropic-rsp-2024 §bio-evals,
citation-pending} or kept lab-internal under
deployment-readiness-review processes
\deepencite{openai-preparedness-2023 §evaluations, citation-pending}.

\subsection{The RSP frame: capability evals as deployment gates}
\label{sec:dangerous-capability-rsp}

The 2023--2024 frontier labs converged on a shared structure for
turning capability evals into deployment commitments. Anthropic's
\emph{Responsible Scaling Policy} \citep{anthropic-rsp-2024}
defines an AI Safety Level (ASL) tier framework. ASL-1 covers
no-autonomy / low-risk systems; ASL-2 is the current frontier
baseline (Claude 3 / Claude 4 family) with basic containment and a
list of capability triggers that must \emph{not} be reached; ASL-3
specifies elevated-risk thresholds on cyber, biology, and autonomy
axes, with deployment requiring hardened security, escalation
procedures, and external evaluation; ASL-4 is reserved for
autonomous-AI risk and remains underspecified by design. The
load-bearing structural commitment is that each tier specifies
\emph{capability evaluations} that must be run before crossing into
the tier and \emph{security and containment commitments} that must be
met \emph{at} the tier; deployment is a conjunction over the two%
\footnote{KB note: \texttt{kb/notes/alignment/safety-evaluation.md\#anthropic-responsible-scaling-policy-rsp}.}.

\begin{forumsignal}
The three frontier-lab framework documents are tier-B sources by the
discipline of \cref{sec:knowledge-benchmarks}: vendor policy
publications, not peer-reviewed papers, and not held to the same
external-replication bar. Anthropic's Responsible Scaling Policy
\citep{anthropic-rsp-2024} defines the ASL tier framework cited
above. OpenAI's \emph{Preparedness Framework}
\deepencite{openai-preparedness-2023 §framework, citation-pending}
defines parallel risk tiers across cybersecurity, CBRN, persuasion,
and model-autonomy axes, with a deployment review process gated on
internal capability scorecards. Google DeepMind's \emph{Frontier
Safety Framework}
\deepencite{deepmind-frontier-safety-2024 §framework, citation-pending}
defines parallel \emph{critical capability levels} with
mitigation-tier commitments. The three frameworks share the layered-
tier structure --- capability triggers cross-checked against
mitigation commitments --- and disagree on specific thresholds and
evaluation protocols%
\footnote{KB note: \texttt{kb/notes/alignment/safety-evaluation.md\#anthropic-responsible-scaling-policy-rsp}; cross-comparison sourced from the AI Index 2025.}.
External evaluators --- UK AISI, US AISI, METR, Apollo Research ---
serve as the external-validation half of these frameworks. Pre-
deployment access agreements between vendors and AISIs are themselves
new institutional infrastructure, on the order of a couple of years
old at the time of writing, and not yet stable.
\end{forumsignal}

The structural contribution of the RSP frame is to commit a vendor,
ahead of capability gain, to specific evaluation protocols and
abandonment thresholds. This is a substantial shift in deployment
discipline relative to the 2020--2022 norm of post-hoc safety
sections in tech reports. It also makes capability evals
\emph{load-bearing in a regulatory sense}: a benchmark whose result
gates deployment is no longer a research instrument, it is part of
the deployment artifact. The dual-use, expert-grading, and
threshold-vs-trend problems below are problems precisely because the
RSP frame raises the consequences of getting them wrong.

\subsection{Methodological problem (a): dual-use publication}
\label{sec:dangerous-capability-dual-use}

The first structural distinction from earlier eval layers is that
the eval material itself is dual-use. A bio-uplift benchmark that
specifies, in publishable detail, exactly which synthesis steps a
model improves over a non-augmented attacker is, by construction, a
manual fragment for the same attacker. A cyber-uplift benchmark
that publishes its full task set publishes a curated set of attacker
training data the next-generation model can train on. The dual-use
constraint pushes the publication frontier in two directions
simultaneously: toward \emph{more redaction} of item-level content
in published reports, and toward \emph{more secure transmission} of
unredacted items between vendor and external evaluator. Neither
direction restores the open-publication norm that adversarial-
robustness benchmarks (\cref{sec:adversarial-robustness}) take for
granted.

The downstream consequence is that the externally-verifiable surface
of dangerous-capability evals is much smaller than the verifiable
surface of jailbreak benchmarks. HarmBench \citep{harmbench2024}
ships $510$ behaviors, $18$ attack methods, and $33$ target LLMs as
a public package. The bio and CBRN evaluations cited in vendor RSP
disclosures
\deepencite{anthropic-rsp-2024 §rsp-evaluations, citation-pending}
do not. A reader who wants to verify a vendor's claim that its
model has \emph{not} crossed an ASL-3 bio-uplift trigger does not
have the option of running the eval; the most a reader can do is
trust the vendor's self-report and the vendor's choice of external
evaluator.

\subsection{Methodological problem (b): expert-graded reproduction}
\label{sec:dangerous-capability-experts}

The second structural distinction is that scoring is expensive. A
multiple-choice MMLU item costs essentially zero to grade once the
gold key is built (\cref{sec:knowledge-mmlu}); a SWE-Bench instance
costs the price of running a test suite (\cref{sec:agentic-swebench}).
A bio-uplift task involves expert scientists evaluating whether the
model's output meaningfully advances a synthesis plan over what a
non-augmented attacker could produce; a cyber-uplift CTF task may
require a security professional's judgment on partial credit and
exploit feasibility; a persuasion eval requires controlled
human-subject comparison against a human-persuader baseline. Expert
graders are slow, expensive, and themselves a finite resource pool.
The eval cycle becomes (i) construct items, (ii) elicit model
outputs, (iii) recruit experts, (iv) grade. Steps (iii) and (iv) cost
weeks to months at scale, which means that the cadence at which
dangerous-capability evals can track frontier capability gain is
fundamentally slower than the cadence at which the underlying models
ship.

\begin{intuition}
Reproducibility on a knowledge benchmark is governed by the cost of
running the verifier $V$ on a fixed item; reproducibility on a
dangerous-capability benchmark is governed by the cost of recruiting
the expert who \emph{is} $V$. The same formal frame from
\cref{sec:agentic-benchmarks} applies --- an instance is still
$\mathcal{B} = (s_0, \mathcal{A}, \mathcal{T}, \rho, g, V)$ --- but
$V$ is human-in-the-loop, expensive, and small-$N$ by
construction. A field that has not built mechanical $V$ has not
built a benchmark in the operational sense; it has built an
\emph{evaluation protocol that produces numbers}. The two are not
the same thing, and treating an expert-graded number as if it were a
benchmark score conflates them. (Returning to canonical form: when
$V$ is human, $\mathrm{ASR}$ from \cref{sec:adversarial-robustness}
becomes a sample-mean estimator over a small expert cohort, and its
variance is dominated by inter-rater disagreement, not by the
behavior of the model.)
\end{intuition}

The second-order effect is that external replication is both rare
and informally negotiated. METR autonomy results
\citep{metr-autonomy-2024} are partially externally verifiable
because the task suite is public and the model APIs are accessible;
bio and CBRN results are not, and the field substitutes
\emph{evaluator-as-replicator} (a separate organization runs the
same protocol on the same model under access agreement) for
\emph{user-as-replicator} (anyone runs the eval). The resulting
trust model is closer to clinical-trial regulation than to
benchmark publication.

\subsection{Methodological problem (c): threshold-vs-trend tension}
\label{sec:dangerous-capability-threshold}

The third structural distinction is the tension between the RSP
frame's commitment to \emph{thresholds} and the empirical reality of
\emph{trends}. RSPs commit a vendor to specific actions when a
capability score crosses a specified line: deploy, deploy with
mitigations, do not deploy. Capability scores, however, are smooth in
model scale, training data, and post-training pipeline (compare the
$1.96\% \to 70\%$ SWE-Bench arc in \cref{sec:agentic-swebench} and
the $15\% \to 50\%+$ GAIA arc; both are continuous trajectories with
no obvious knee). Setting a threshold on a smooth curve forces a
discrete decision out of a continuous quantity, and small differences
in score near the threshold produce large differences in deployment
posture. Two failure modes follow.

The first is \emph{threshold sandbagging}: incentives push toward
designing or interpreting the eval such that the model's score lands
just below the threshold, since crossing it imposes substantial
mitigation cost. This is structurally analogous to the
sycophancy-in-rater-data Goodhart pattern of
\cref{sec:sycophancy-rlhf}: when the eval is itself a
preference-ranked artifact (vendor design vs.\ vendor commitment),
the gradient on item difficulty and grading rubric is not aligned
with the gradient on safety. The second is \emph{threshold
discontinuity surprise}: when a smoothly-improving capability does
cross the line, the discrete commitment kicks in and the deployment
posture changes abruptly, which the operational organization may not
be prepared to absorb. Smooth capability gain plus discrete policy
response is not a stable engineering combination.

The interaction with the dual-use and expert-grading problems
compounds. If the eval is dual-use and expert-graded, the surface
on which a vendor and an external evaluator can disagree about
whether the threshold has been crossed is large, and the ability of
the public to adjudicate is small. The discrete commitment is built
on top of a small-$N$, expert-judged, partially-redacted measurement
--- which is not the substrate one would choose for high-confidence
deployment decisions if the measurement substrate could be designed
freely.

\begin{contradiction}
\textbf{Dangerous-capability eval reproducibility.} The field's
deployment safety story for the bio, cyber, persuasion, and autonomy
axes rests on capability evaluations whose item content is partially
redacted (dual-use), whose grading is expert-mediated (expensive
and small-$N$), and whose threshold commitments are pegged to
smoothly-trending capability scores. The result is that external
parties cannot openly verify vendor claims that a given model is, or
is not, above a given RSP-defined trigger. The vendor side argues
that this regime is the best feasible given dual-use constraints
and that external-evaluator access (UK AISI, US AISI, METR, Apollo)
substitutes for open verification. The critic side argues that
deployment-gating decisions of this consequence cannot rest on a
trust model the public cannot audit, and that the absence of a
field-standard reproducibility protocol is itself a methodological
failure. Both readings are defensible from the available evidence as
of 2026-Q2. The contradiction would close along one of two paths:
(i) standardized, publicly-audited evaluator access agreements with
post-hoc disclosure of methodology and aggregate results, or (ii) a
class of dangerous-capability evals constructed such that item-level
redaction is unnecessary because the capability of interest can be
measured on a non-dual-use proxy. Neither path has shipped at field
scale.
\end{contradiction}

\subsection{Forward-link: the alignment-threat layer}
\label{sec:dangerous-capability-forward}

The threats taken up in
\cref{sec:sycophancy,sec:scheming,sec:alignment-faking} shift the
construct of evaluation in a different direction. Where this section
asks whether the model has dangerous capabilities, the next layer
asks whether the model's \emph{behavior under deployment} aligns
with its operator's intent. Sycophancy is a Goodhart artifact of
preference-RL that does not require capability for deception
(\cref{sec:sycophancy}). Scheming is capability for in-context
strategic deception under an instructed goal
(\cref{sec:scheming}). Alignment-faking is strategy under training
distribution shift, and reuses the legible-CoT signal as part of its
detection surface (\cref{sec:alignment-faking}). The cross-paper
thread runs back to the post-training pipeline of Paper 2: RSPs
intervene at the SFT and RL boundary --- the point at which a
trained model becomes a deployable artifact --- gating before
deployment rather than re-training after the fact. This locates
dangerous-capability evals at exactly the post-training-to-deployment
transition where the alignment-threat layer also lives.
\Cref{sec:scalable-oversight} closes the layered defense by asking
how any of this generalizes when capability passes the level at which
the human evaluator can directly verify the model's outputs.
